% !TeX root = mainstarkov.tex
% !TeX program = xelatex
% !TeX encoding = UTF-8
\newpage
\section*{РЕЗЮМЕ СТАРТАП-ПРОЕКТА}
\addcontentsline{toc}{section}{РЕЗЮМЕ СТАРТАП-ПРОЕКТА}

Название: «\Тема\ \ТемаВтораяСтрока\ \ТемаТретьяСтрока». Проект представляет собой кроссплатформенное приложение IndoorNav, предназначенное для навигации внутри зданий и автоматического построения эвакуационных маршрутов в чрезвычайных ситуациях. Система обеспечивает построение кратчайших маршрутов между помещениями, поиск ближайшего эвакуационного выхода, динамическое перестроение маршрута при блокировке участков пути, а также реализацию комплекса мер по обеспечению безопасности пользователей в режиме ЧС.

Ключевые возможности системы включают:
\begin{enumerate}
\item Построение кратчайших маршрутов между произвольными точками внутри здания с учётом межэтажных переходов и длин рёбер навигационного графа.

\item Автоматический поиск ближайшего эвакуационного выхода при активации режима чрезвычайной ситуации.

\item Перестроение маршрута в обход заблокированных участков без повторной загрузки графа.

\item Определение стартовой точки пользователя путём сканирования QR-кодов со стены здания.

\item Определение аудитории студента по расписанию и построение пути к выходу после её подтверждения.

\item Встроенный редактор графа, позволяющий администратору формировать карту здания без перекомпиляции приложения.
\end{enumerate}

Разработанная система оптимизирует процесс эвакуации из замкнутых пространств, повышая безопасность пользователей за счёт оперативного построения оптимальных маршрутов к выходам с учётом заблокированных проходов.

Актуальность проекта обусловлена требованиями обеспечения безопасности людей в зданиях с большим количеством помещений и сложной планировкой. Бумажный план эвакуации статичен и не учитывает блокировку отдельных участков, что в условиях чрезвычайной ситуации может привести к выбору неоптимального или недоступного маршрута — цифровая система подобной ошибки не допускает.

Целью разработки является создание системы построения эвакуационных маршрутов в реальном времени, обеспечивающей:
\begin{itemize}
\item оперативное построение кратчайших маршрутов до ближайших выходов с учётом текущей обстановки;
\item динамическое изменение маршрута при блокировке участков пути;
\item работу на мобильных устройствах и настольных компьютерах;
\item возможность масштабирования на здания с различной планировкой и этажности.
\end{itemize}

Отличительные особенности разработанной системы по сравнению с существующими решениями:
\begin{enumerate}
\item Алгоритм Дейкстры с возможностью динамического исключения узлов, позволяющий перестраивать маршрут без модификации исходного графа.

\item Кроссплатформенная реализация на базе .NET MAUI с поддержкой Android, iOS и Windows.

\item Полная автономность: данные хранятся локально, потеря сети не прерывает навигацию.

\item Редактор карты для администратора, позволяющий формировать навигационный граф поверх плана этажа.

\item Интеграция с расписанием занятий для автоматического определения местоположения студента при активации ЧС.
\end{enumerate}

В системе реализованы алгоритмы маршрутизации (Дейкстра с исключением узлов, поиск ближайшего выхода), сервис управления режимом ЧС, аутентификация с разделением ролей (администратор, студент, гость), управление расписанием и позиционирование по QR-кодам.

Идея создания проекта возникла в 2025 году в связи с необходимостью обеспечения безопасной и оперативной эвакуации в зданиях университета. Разработка системы была начата в 2025 году и завершена в 2026 году в рамках выполнения выпускной квалификационной работы.

Бизнес-цели:
\begin{enumerate}
\item Внедрение системы навигации в зданиях университета для повышения безопасности студентов и сотрудников при чрезвычайных ситуациях.

\item Масштабирование системы для использования в зданиях различного назначения: образовательные учреждения, бизнес-центры, торговые комплексы.

\item Интеграция с системами пожарной безопасности и оповещения для автоматической активации режима ЧС.

\item Добавление push-уведомлений о чрезвычайных ситуациях и голосовые подсказки при навигации.

\item Ускорение процесса эвакуации за счёт автоматического определения местоположения пользователя.
\end{enumerate}
